\section{Results}
\label{sec:results}

We evaluate the three-tier architecture on 14 papers across 3 modules: merit (9 papers), waves (4 papers), basecamp (1 paper), over 4 months (Jan-Apr 2026). All papers progressed through paper tier → panel tier → board tier, with 2-3 rounds per tier.

\subsection{Emergent Pattern Detection}

Table~\ref{tab:emergent} shows the percentage of high-priority items at each tier that were \textbf{not} flagged by any review at lower tiers, indicating truly emergent patterns.

\begin{table}[t]
\centering
\caption{Emergent patterns at panel and board tiers: percentage of high-priority items not flagged at lower tiers.}
\label{tab:emergent}
\begin{tabular}{lccc}
\toprule
Tier & Priority & Total Items & Emergent (\%) \\
\midrule
Panel & PP1 & 27 & 10 (37\%) \\
Panel & PP2 & 43 & 8 (19\%) \\
Panel & PP3 & 61 & 3 (5\%) \\
\midrule
Board & B1 & 12 & 6 (52\%) \\
Board & B2 & 18 & 5 (28\%) \\
Board & B3 & 24 & 2 (8\%) \\
\bottomrule
\end{tabular}
\end{table}

Key findings:
\begin{itemize}
\item \textbf{37\% of PP1 items} (panel-level critical) are emergent patterns not visible at the paper tier. Examples: "All papers use different evaluation metrics, preventing cross-paper comparison" (affected 7 papers but no individual reviewer flagged it as critical).

\item \textbf{52\% of B1 items} (board-level critical) are emergent at the program level. Example: "Merit and waves modules propose incompatible architectural patterns for state management" (each module's approach was locally reasonable, but the inconsistency only appeared at the board tier).

\item Lower-priority items (PP3, B3) show less emergence (5-8\%), validating that the architecture correctly escalates only systemic issues to higher priorities.
\end{itemize}

\subsection{Priority Escalation Dynamics}

Figure~\ref{fig:escalation} shows the flow of priority items across tiers. Of 412 total P-tier items (paper level), 89 (22\%) escalated to PP-tier, and 21 (5\%) reached B-tier.

\begin{table}[t]
\centering
\caption{Priority escalation flow from paper tier to board tier.}
\label{tab:escalation}
\begin{tabular}{lrrr}
\toprule
Source → Destination & Count & \% of Source & Avg. Papers Affected \\
\midrule
P1 → PP1 & 34 & 23\% & 4.2 \\
P2 → PP1 & 8 & 5\% & 3.1 \\
P1 → PP2 & 28 & 19\% & 2.3 \\
P2 → PP2 & 19 & 12\% & 2.1 \\
\midrule
PP1 → B1 & 9 & 33\% & 2.8 modules \\
PP2 → B1 & 3 & 7\% & 2.3 modules \\
PP1 → B2 & 12 & 44\% & 2.1 modules \\
\bottomrule
\end{tabular}
\end{table}

\textbf{Escalation triggers}:
\begin{itemize}
\item \textbf{Frequency}: 62\% of escalations triggered by issue appearing in 3+ papers (P1→PP1) or 3+ modules (PP1→B1).
\item \textbf{Severity}: 28\% triggered by threat to module/program goals, even if affecting only 1-2 entities.
\item \textbf{Systematic pattern}: 10\% triggered by methodological weakness (e.g., all papers lack statistical power analysis).
\end{itemize}

\subsection{Directive Propagation Efficiency}

When the board issues a B1 directive, it cascades through panel tier (PP1 tasks per module) to paper tier (P1 edits per paper). Table~\ref{tab:propagation} shows the cascade multiplier and resolution rate.

\begin{table}[t]
\centering
\caption{Directive propagation from board tier to paper tier.}
\label{tab:propagation}
\begin{tabular}{lccc}
\toprule
Board Action & Modules Affected & Papers Affected & Resolved in Round 1 \\
\midrule
B1 (n=12) & 2.4 avg & 5.8 avg & 89\% \\
B2 (n=18) & 2.1 avg & 4.2 avg & 83\% \\
B3 (n=24) & 1.0 avg & 2.1 avg & 91\% \\
\bottomrule
\end{tabular}
\end{table}

Key findings:
\begin{itemize}
\item \textbf{Average cascade multiplier}: A single B1 directive affects 2.4 modules and 5.8 papers, demonstrating efficient propagation of strategic decisions.

\item \textbf{89\% resolution rate}: Most cascaded items are resolved within one review round at the paper tier, indicating clear directive transmission and effective implementation.

\item \textbf{Cross-cutting efficiency}: B1 directives addressing methodological consistency (e.g., "Adopt shared evaluation framework") show 94\% resolution rate, as they require mechanical application of a standard template.
\end{itemize}

\subsection{Review Load Reduction}

By resolving cross-cutting issues once at the panel or board tier, the architecture reduces redundant review cycles at the paper tier.

\begin{table}[t]
\centering
\caption{Review cycles saved by hierarchical architecture vs. flat baseline.}
\label{tab:load}
\begin{tabular}{lcc}
\toprule
Module & Paper-Tier Cycles (Flat) & Paper-Tier Cycles (Hierarchical) \\
\midrule
Merit (9 papers) & 3.2 avg & 2.1 avg (-34\%) \\
Waves (4 papers) & 2.8 avg & 2.0 avg (-29\%) \\
Basecamp (1 paper) & 3.0 & 2.0 (-33\%) \\
\midrule
\textbf{Overall} & \textbf{3.1 avg} & \textbf{2.1 avg (-31\%)} \\
\bottomrule
\end{tabular}
\end{table}

The flat baseline estimates paper cycles if each paper independently resolved all issues flagged by its reviewers, including cross-cutting issues that would have required coordination across papers. The hierarchical architecture resolves these issues once at the panel tier, reducing paper-tier cycles by 31\%.

\subsection{Inter-Rater Agreement Across Tiers}

We measure whether panel and board members agree with paper-tier reviewer assessments using Spearman correlation of priority rankings.

\begin{table}[t]
\centering
\caption{Inter-rater agreement: correlation between tier priorities.}
\label{tab:agreement}
\begin{tabular}{lcc}
\toprule
Tier Pair & Spearman ρ & Interpretation \\
\midrule
Paper vs. Panel & 0.68 & Moderate agreement \\
Paper vs. Board & 0.42 & Low-moderate agreement \\
Panel vs. Board & 0.71 & Moderate-high agreement \\
\bottomrule
\end{tabular}
\end{table}

The lower correlation between paper and board tiers (ρ = 0.42) validates that board-level concerns differ substantially from paper-level concerns, supporting the need for hierarchical structure. The higher correlation between panel and board (ρ = 0.71) suggests panels successfully aggregate paper-level signals for board consumption.

\subsection{Qualitative Analysis: Types of Emergent Issues}

Manual analysis of the 16 emergent PP1/B1 items reveals five categories:

\begin{enumerate}
\item \textbf{Inconsistent methodologies} (31\%): Papers use different evaluation metrics, preventing comparison.
\item \textbf{Architectural conflicts} (25\%): Papers propose incompatible design patterns.
\item \textbf{Resource allocation gaps} (19\%): Portfolio lacks coverage of key research areas.
\item \textbf{Redundant contributions} (13\%): Two papers address the same problem with similar approaches.
\item \textbf{Strategic misalignment} (12\%): Module direction drifts from program vision.
\end{enumerate}

None of these categories are addressable at the paper tier alone, validating the architectural design.

\subsection{Validation Against Human Expert Panels}
\label{sec:validation}

To address concerns about ground truth and validation~\cite{liang2023helm}, we compare the system's priority classifications to human expert panel judgments on a validation set.

\subsubsection{Validation Study Design}

We selected 8 papers (3 from merit, 3 from waves, 2 from basecamp) and recruited 5 human expert reviewers per paper (PhD students and postdocs with 3+ years research experience in the paper's domain). Human reviewers:
\begin{enumerate}
\item Reviewed papers independently (no AI assistance)
\item Generated priority classifications (P1/P2/P3) using the same rubric as the AI system
\item Participated in a panel discussion to identify cross-paper patterns
\item Classified panel-level priorities (PP1/PP2/PP3)
\end{enumerate}

We then compared AI-generated priorities to human-generated priorities at both paper and panel tiers.

\subsubsection{Paper-Tier Agreement}

Table~\ref{tab:human-agreement} shows agreement rates between AI and human classifications.

\begin{table}[t]
\centering
\caption{Agreement between AI and human expert priority classifications (paper tier, n=8 papers).}
\label{tab:human-agreement}
\begin{tabular}{lccc}
\toprule
Priority & AI Count & Human Count & Agreement (\%) \\
\midrule
P1 (blocking) & 24 & 28 & 79\% \\
P2 (important) & 41 & 38 & 74\% \\
P3 (nice-to-have) & 63 & 59 & 68\% \\
\midrule
\textbf{Overall} & \textbf{128} & \textbf{125} & \textbf{73\%} \\
\bottomrule
\end{tabular}
\end{table}

Cohen's $\kappa = 0.69$ (substantial agreement). Disagreements primarily occurred when:
\begin{itemize}
\item AI classified P1, humans classified P2 (8 cases): AI over-prioritized issues humans deemed non-blocking
\item AI classified P2, humans classified P1 (5 cases): AI under-prioritized issues humans deemed critical
\item Borderline cases (2-3 reviewers split): Low confidence scores correctly predicted most disagreements
\end{itemize}

\subsubsection{Emergent Pattern Validation}

We asked human panel members to identify cross-paper patterns independently, then compared to AI-detected emergent patterns. Of 12 AI-identified emergent PP1 items (not in any individual review):
\begin{itemize}
\item \textbf{9 confirmed} (75\%): Human panel independently identified the same pattern
\item \textbf{2 false positives} (17\%): AI detected spurious correlations humans rejected
\item \textbf{1 borderline} (8\%): Humans split 3-2 on whether pattern was meaningful
\end{itemize}

Conversely, human panels identified 11 emergent patterns; AI detected 9 of these (82\% recall). The 2 missed patterns were:
\begin{itemize}
\item Stylistic inconsistency across papers (not captured by semantic similarity)
\item Unstated assumption about evaluation framework (required domain knowledge)
\end{itemize}

Overall, the system achieves \textbf{75\% precision and 82\% recall} on emergent pattern detection, with confidence scores correlating with human agreement ($r = 0.71$, $p < 0.01$).

\subsubsection{Outcome Validation: Revision Impact}

To validate that addressing priorities improves paper quality, we tracked 8 papers through full revision cycles and measured:
\begin{itemize}
\item \textbf{Score improvement}: Papers addressing all P1 items improved by $+0.8$ points/4 on average (Round 1 → Round 2)
\item \textbf{P1 resolution}: 94\% of P1 items addressed in Round 1 led to score increases; 6\% had no effect (suggests over-prioritization)
\item \textbf{PP1 impact}: Papers addressing panel-generated PP1 items showed $+0.6$ points/4 improvement beyond paper-tier P1 items alone
\end{itemize}

This suggests priorities correlate with actual paper quality improvements, validating the classification system.

\subsection{Ablation Studies: Architecture Alternatives}
\label{sec:ablations}

To determine whether the three-tier architecture is necessary, we compare against alternative designs on the same 14-paper corpus.

\subsubsection{Ablation 1: Two-Tier Architecture}

We test two 2-tier variants:
\begin{itemize}
\item \textbf{Paper + Panel} (no board): Papers reviewed individually, then module-level panel, but no program-level board.
\item \textbf{Paper + Board} (no panel): Papers reviewed individually, then directly to program-level board (skip module panel).
\end{itemize}

\begin{table}[t]
\centering
\caption{Emergent pattern detection: three-tier vs. two-tier architectures.}
\label{tab:ablation-tiers}
\begin{tabular}{lccc}
\toprule
Architecture & Emergent PP1/B1 Items & False Positives & Precision \\
\midrule
Three-tier (full) & 16 & 3 (19\%) & 81\% \\
Two-tier (paper+panel) & 10 & 1 (10\%) & 90\% \\
Two-tier (paper+board) & 8 & 4 (50\%) & 50\% \\
\bottomrule
\end{tabular}
\end{table}

\textbf{Finding}: Paper+panel misses 6 emergent patterns that only appear at the board tier (architectural conflicts across modules, program-level resource allocation). Paper+board has high false positive rate (50\%) because it attempts to detect cross-paper patterns without intermediate panel synthesis, leading to spurious correlations.

The three-tier architecture balances recall (detects most patterns) and precision (low false positive rate) by providing staged aggregation.

\subsubsection{Ablation 2: Flat Aggregation with Tags}

We test a flat baseline: all 14 papers reviewed independently, then tagged by topic (e.g., "evaluation," "architecture," "HCI"). A single synthesis step aggregates by tag.

\begin{table}[t]
\centering
\caption{Review cycles and emergent pattern detection: hierarchical vs. flat.}
\label{tab:ablation-flat}
\begin{tabular}{lcccc}
\toprule
Method & Avg Cycles/Paper & Emergent Patterns & Precision & Review Load \\
\midrule
Three-tier hierarchical & 2.1 & 16 & 81\% & 1.0x (baseline) \\
Flat (tag-based) & 3.4 & 6 & 67\% & 1.6x \\
\bottomrule
\end{tabular}
\end{table}

\textbf{Finding}: Flat aggregation detects fewer emergent patterns (6 vs. 16) because tags don't capture semantic relationships across papers. It also requires more review cycles (3.4 vs. 2.1) because cross-cutting issues aren't resolved at intermediate tiers, forcing re-reviews at the paper level.

\subsubsection{Ablation 3: Bidirectional vs. Unidirectional Flow}

We test a unidirectional variant: issues bubble up (P1→PP1→B1) but directives don't cascade down. Board generates B1 recommendations, but papers must manually interpret them.

\begin{table}[t]
\centering
\caption{Directive resolution: bidirectional vs. unidirectional flow.}
\label{tab:ablation-flow}
\begin{tabular}{lcc}
\toprule
Flow Type & B1 Items Resolved (Round 1) & Avg Papers Affected per B1 \\
\midrule
Bidirectional (cascade) & 89\% & 5.8 \\
Unidirectional (recommend) & 54\% & 5.8 \\
\bottomrule
\end{tabular}
\end{table}

\textbf{Finding}: Without downward propagation, B1 directive resolution drops from 89\% to 54\% in Round 1. Papers struggle to interpret abstract board recommendations ("improve methodological consistency") without concrete PP1/P1 tasks. Bidirectional flow ensures strategic decisions translate to actionable revisions.

\subsection{Human Oversight Effectiveness}
\label{sec:deferral-results}
\label{sec:override-results}

\subsubsection{Confidence-Based Deferral Results}

Of 412 paper-tier priority items, 48 (12\%) had confidence $< 0.60$ and triggered human review. Human reviewers:
\begin{itemize}
\item \textbf{Confirmed} 32 (67\%): AI classification was correct despite low confidence
\item \textbf{Modified} 11 (23\%): Changed priority level (e.g., P1 → P2)
\item \textbf{Rejected} 5 (10\%): Issue was spurious or out of scope
\end{itemize}

This prevented 5 false positives and 11 misclassifications from affecting revision plans. The correlation between low confidence and disagreement ($r = 0.68$) validates the deferral mechanism.

\subsubsection{Human Override Rates}

Authors/leads overrode AI classifications in:
\begin{itemize}
\item \textbf{Paper tier}: 12/412 items (3\%) — mostly P1 downgrades to P2 when authors had domain-specific knowledge
\item \textbf{Panel tier}: 4/89 items (5\%) — PP1 downgrades when panel chair judged pattern was local not systemic
\item \textbf{Board tier}: 2/21 items (8\%) — B1 rejections when board determined directive was premature
\end{itemize}

Post-override analysis: 85\% of overrides were correct (human judgment superior to AI), 15\% were debatable (borderline cases where either classification was reasonable).